\documentclass[11pt]{article}

\usepackage[a4paper,margin=1in]{geometry}
\usepackage{amsmath,amssymb,amsthm,mathtools}
\usepackage{booktabs}
\usepackage{hyperref}

\newtheorem{theorem}{Theorem}[section]
\newtheorem{lemma}[theorem]{Lemma}
\newtheorem{proposition}[theorem]{Proposition}
\newtheorem{corollary}[theorem]{Corollary}
\theoremstyle{definition}
\newtheorem{definition}[theorem]{Definition}
\newtheorem{remark}[theorem]{Remark}

\DeclareMathOperator{\Asym}{\mathcal{A}}
\newcommand{\bdy}{\partial}
\newcommand{\Rn}{\mathbb{R}^N}
\newcommand{\Ot}{\widetilde{\Omega}}
\newcommand{\HN}{\mathcal{H}^{N-1}}
\newcommand{\NS}{\mathrm{NearlySphericalClosure}}

\title{Nearly Spherical Closure for Route~A:\\
sharp Saint--Venant stability via BDV Theorem~3.3 and Lemma~4.2}
\author{Plan~1 / Route~A --- block \texttt{NearlySphericalClosure}}
\date{}

\begin{document}
\maketitle

\begin{abstract}
We record, as a self-contained pointwise quantitative statement, the
``nearly spherical closure'' building block of Route~A. Starting from two
quantitative results in Brasco--De~Philippis--Velichkov~\cite{BDV2013}
(Theorem~3.3 and Lemma~4.2), we produce, with explicit dimensional
constants, the one-step inequality
\[
E(\Ot)-E(B_1)\;\ge\; c_*(N)\,\alpha(\Ot)
\quad\text{and}\quad
E(\Ot)-E(B_1)\;\ge\;c_*(N)\,|B_1|^{2}\,C_1(N)^{-1}\,\Asym(\Ot)^{2},
\]
valid for every $C^{2,\gamma_0}$ nearly spherical set $\Ot$ with
sufficiently small parametrisation. Here $E$ denotes the Saint--Venant
(torsion) energy used by BDV; the transfer to the first Dirichlet
eigenvalue is performed downstream by the \texttt{KohlerJobinTransfer}
block. The whole argument is by direct substitution into the BDV
estimates: no compactness, no contradiction and no selection principle
are used here.
\end{abstract}

\section{Setup, conventions and the BDV estimates}\label{sec:setup}

Throughout, $N\ge 2$ and $0<\gamma_0\le 1$ are fixed. We write
$|\cdot|=\mathcal{L}^N$ for Lebesgue measure and $\HN$ for the surface
measure on $\bdy B_1$.

\paragraph{Saint--Venant (torsion) energy.}
Following \cite[\S2]{BDV2013}, for every open bounded set
$\Omega\subset\Rn$ we denote by $u_\Omega\in W^{1,2}_0(\Omega)$ the unique
weak solution of
\[
-\Delta u_\Omega=1\quad\text{in }\Omega,\qquad u_\Omega=0\quad\text{on }
\bdy\Omega,
\]
and by
\[
T(\Omega)\;:=\;\int_\Omega u_\Omega\,dx
\]
the torsional rigidity. The Saint--Venant energy is
\begin{equation}\label{eq:Edef}
E(\Omega)\;:=\;-\tfrac12\,T(\Omega)
\;=\;\inf_{v\in W^{1,2}_0(\Omega)}
\Bigl(\tfrac12\!\int_\Omega|\nabla v|^2\,dx-\int_\Omega v\,dx\Bigr).
\end{equation}
This is the energy stabilised by Theorem~3.3 of \cite{BDV2013}. The
transfer to the first Dirichlet eigenvalue
$\lambda_1(\Omega)=\inf\{\int|\nabla u|^2:u\in W^{1,2}_0(\Omega),
\|u\|_{L^2}=1\}$ is performed by the K\"ohler--Jobin inequality and is
the subject of a separate block (\texttt{KohlerJobinTransfer}).

\paragraph{Fraenkel and smooth asymmetries.}
The (classical) \emph{Fraenkel asymmetry} is
\[
\Asym(\Omega)\;:=\;\inf_{x_0\in\Rn}
\frac{|\Omega\triangle B_1(x_0)|}{|B_1|},
\qquad |\Omega|=|B_1|.
\]
Following \cite[Definition~4.1]{BDV2013} we also introduce the
\emph{smooth asymmetry}
\begin{equation}\label{eq:alphadef}
\alpha(\Omega)\;:=\;\int_{\Omega\triangle B_1(x_\Omega)}
\bigl|\,1-|x-x_\Omega|\,\bigr|\,dx,
\qquad x_\Omega=\tfrac1{|\Omega|}\!\int_\Omega x\,dx.
\end{equation}
The smooth asymmetry $\alpha$ weights the symmetric difference by the
distance to the unit sphere, while $\Asym$ uses the unweighted symmetric
difference. The inequality between them is Lemma~\ref{lem:BDV42}(i)
below; it goes in the direction $\Asym^{2}\lesssim\alpha$.

\paragraph{Nearly spherical sets.}
Following \cite[Definition~3.1]{BDV2013} we say that an open bounded set
$\Ot\subset\Rn$ is \emph{nearly spherical of class $C^{2,\gamma_0}$ with
parametrisation $g$} if there exists $g\in C^{2,\gamma_0}(\bdy B_1)$
with $\|g\|_{L^\infty(\bdy B_1)}\le \tfrac12$ such that
\[
\bdy\Ot=\bigl\{(1+g(y))\,y:y\in\bdy B_1\bigr\},\qquad
\Ot=\bigl\{r\,y:y\in\bdy B_1,\ 0\le r<1+g(y)\bigr\}.
\]
The $H^{1/2}(\bdy B_1)$ norm is, as in \cite[Definition~3.2]{BDV2013},
\[
\|g\|_{H^{1/2}(\bdy B_1)}^2
\;:=\;\int_{\bdy B_1}g^2\,d\HN
\;+\;\int_{B_1}|\nabla H(g)|^2\,dx,
\]
where $H(g)$ is the $W^{1,2}$-harmonic extension of $g$ to $B_1$.

\medskip
The two quantitative statements we use as a black box are now recalled.

\begin{lemma}[BDV Theorem~3.3, sharp nearly spherical stability]
\label{lem:BDV33}
For every $0<\gamma_0\le 1$ there exists $\delta_1=\delta_1(N,\gamma_0)>0$
such that, for every nearly spherical set $\Omega$ of class
$C^{2,\gamma_0}$ parametrised by $g$ with
\[
\|g\|_{C^{2,\gamma_0}(\bdy B_1)}\le\delta_1,\qquad
|\Omega|=|B_1|,\qquad x_\Omega=0,
\]
one has
\begin{equation}\label{eq:BDV33}
E(\Omega)-E(B_1)\;\ge\;\frac{1}{32\,N^2}\,\|g\|_{H^{1/2}(\bdy B_1)}^2,
\end{equation}
with $E$ the Saint--Venant energy \eqref{eq:Edef}.
\end{lemma}

The dimensional constant
$c_{\rm sph}(N):=\tfrac{1}{32\,N^2}$ is the value of the bilinear-form
coefficient appearing in \cite[(3.9)--(3.10)]{BDV2013}; it comes from the
gap to the second Steklov eigenvalue of $B_1$ (equal to $2$ in every
dimension).

\begin{lemma}[BDV Lemma~4.2]\label{lem:BDV42}
There exist constants $C_1=C_1(N)>0$, $\delta_3=\delta_3(N)>0$ and
$C_3=C_3(N)>0$ with the following properties.
\begin{enumerate}
\item[\textup{(i)}] For every measurable set $\Omega\subset\Rn$ with
$|\Omega|=|B_1|$,
\begin{equation}\label{eq:BDV42i}
C_1(N)\,\alpha(\Omega)\;\ge\;|\Omega\triangle B_1(x_\Omega)|^2
\;\ge\;|B_1|^2\,\Asym(\Omega)^2.
\end{equation}
\item[\textup{(iii)}] For every nearly spherical set $\Omega$
parametrised by $g$ with $\|g\|_{L^\infty(\bdy B_1)}\le\delta_3$,
\begin{equation}\label{eq:BDV42iii}
\alpha(\Omega)\;\le\;C_3(N)\,\|g\|_{L^2(\bdy B_1)}^2.
\end{equation}
\end{enumerate}
\end{lemma}

\begin{proof}[Reference]
Items (i) and (iii) are \cite[Lemma~4.2~(i),(iii)]{BDV2013}; the proof
there is by direct rearrangement / Taylor expansion of the weighted
integral~\eqref{eq:alphadef}. We use only these two items in the present
section; in particular we do \emph{not} use the BDV selection principle
nor Lemma~4.2(ii).
\end{proof}

\section{The nearly spherical closure}\label{sec:closure}

\subsection{Statement in the volume- and barycenter-normalised case}

The trivial trace--Poincar\'e comparison
\begin{equation}\label{eq:Hhalfvsi2}
\|g\|_{L^2(\bdy B_1)}^2 \;\le\; \|g\|_{H^{1/2}(\bdy B_1)}^2,
\end{equation}
follows directly from the definition of $\|\cdot\|_{H^{1/2}(\bdy B_1)}$
($L^2$ is one of the two summands).

\begin{theorem}[$\NS$: nearly spherical closure]\label{thm:NS}
Fix $0<\gamma_0\le 1$ and set
\begin{equation}\label{eq:dsph}
\delta_{\rm sph}(N,\gamma_0)
\;:=\;\min\!\bigl(\delta_1(N,\gamma_0),\,\delta_3(N)\bigr).
\end{equation}
For every $C^{2,\gamma_0}$ nearly spherical set
$\Ot=\{(1+g(y))y:y\in\bdy B_1\}$ with
\begin{equation}\label{eq:normal}
\|g\|_{C^{2,\gamma_0}(\bdy B_1)}\le\delta_{\rm sph}(N,\gamma_0),
\qquad |\Ot|=|B_1|,\qquad x_{\Ot}=0,
\end{equation}
we have
\begin{align}
E(\Ot)-E(B_1)
&\;\ge\;
\frac{c_{\rm sph}(N)}{C_3(N)}\,\alpha(\Ot)
\;=:\;c_*(N)\,\alpha(\Ot),
\label{eq:NSalpha}\\
E(\Ot)-E(B_1)
&\;\ge\;
\frac{c_{\rm sph}(N)\,|B_1|^{2}}{C_3(N)\,C_1(N)}\,\Asym(\Ot)^{2}
\;=\;c_*(N)\,|B_1|^{2}\,C_1(N)^{-1}\,\Asym(\Ot)^{2},
\label{eq:NSfraenkel}
\end{align}
with the explicit dimensional constant
\begin{equation}\label{eq:cstar}
c_*(N)\;=\;\frac{c_{\rm sph}(N)}{C_3(N)}
\;=\;\frac{1}{32\,N^2\,C_3(N)}.
\end{equation}
Equivalently, the smooth-asymmetry energy quotient is bounded below by an
explicit dimensional constant:
\begin{equation}\label{eq:Qalpha}
Q_\alpha(\Ot)\;:=\;\frac{E(\Ot)-E(B_1)}{\alpha(\Ot)}
\;\ge\;c_*(N).
\end{equation}
\end{theorem}

\begin{proof}
By~\eqref{eq:dsph}, the hypotheses of Lemma~\ref{lem:BDV33}
(threshold $\delta_1$) and of Lemma~\ref{lem:BDV42}(iii)
(threshold $\delta_3$) hold simultaneously.

\medskip\noindent\emph{Step 1: smooth-asymmetry form~\eqref{eq:NSalpha}.}\;
Chain Lemma~\ref{lem:BDV33}, the comparison~\eqref{eq:Hhalfvsi2}, and
Lemma~\ref{lem:BDV42}(iii):
\[
E(\Ot)-E(B_1)
\stackrel{\eqref{eq:BDV33}}{\ge}
\tfrac{1}{32N^2}\,\|g\|_{H^{1/2}(\bdy B_1)}^2
\stackrel{\eqref{eq:Hhalfvsi2}}{\ge}
\tfrac{1}{32N^2}\,\|g\|_{L^2(\bdy B_1)}^2
\stackrel{\eqref{eq:BDV42iii}}{\ge}
\tfrac{1}{32N^2\,C_3(N)}\,\alpha(\Ot).
\]
This is~\eqref{eq:NSalpha}.

\medskip\noindent\emph{Step 2: Fraenkel form~\eqref{eq:NSfraenkel}.}\;
Combining with Lemma~\ref{lem:BDV42}(i):
\[
E(\Ot)-E(B_1)
\;\ge\;c_*(N)\,\alpha(\Ot)
\stackrel{\eqref{eq:BDV42i}}{\ge}
\frac{c_*(N)}{C_1(N)}\,|B_1|^{2}\,\Asym(\Ot)^{2}.
\]

The quotient identity~\eqref{eq:Qalpha} is~\eqref{eq:NSalpha} rewritten.
\end{proof}

\begin{remark}[Direction of asymmetry inequalities]
The inequalities~\eqref{eq:BDV42i} and~\eqref{eq:BDV42iii} go in opposite
directions: $\alpha$ is \emph{lower bounded} by $\Asym^{2}$ (so $\alpha$
is the larger, easier-to-bound-below quantity) and \emph{upper bounded}
by $\|g\|_{L^2(\bdy B_1)}^2$. Theorem~\ref{thm:NS} uses both in the
favourable direction: \eqref{eq:BDV42iii} converts the $H^{1/2}$ energy
lower bound into an $\alpha$-lower bound, while~\eqref{eq:BDV42i}
transfers it to $\Asym^{2}$.
\end{remark}

\subsection{Absorbing the volume/barycenter normalisation}

In Route~A the set $\Ot$ produced upstream is $C^{2,\gamma_0}$-close to
$B_1$ but may have $|\Ot|\neq|B_1|$ or $x_{\Ot}\neq 0$. The normalisation
is contained in BDV's own proof of Theorem~3.3 (see
\cite[(3.6)--(3.7)]{BDV2013}): expanding the two scalar identities
\begin{equation}\label{eq:vol}
|\Ot|=\int_{\bdy B_1}\frac{(1+g)^N}{N}\,d\HN=|B_1|,
\end{equation}
\begin{equation}\label{eq:bary}
x_{\Ot}=\int_{\bdy B_1}y\,\frac{(1+g)^{N+1}}{N+1}\,d\HN=0,
\end{equation}
in $g$ and using $\|g\|_{C^{2,\gamma_0}}\le\delta_1$ yields the moment
bounds \cite[(3.6)--(3.7)]{BDV2013}
\begin{equation}\label{eq:moment}
\Bigl|\int_{\bdy B_1}g\,d\HN\Bigr|
+\sum_{i=1}^{N}\Bigl|\int_{\bdy B_1}y_i\,g\,d\HN\Bigr|
\;\le\;C(N)\,\delta_1\,\|g\|_{H^{1/2}(\bdy B_1)}.
\end{equation}

\begin{corollary}[Normalisation absorbed in the constant]
\label{cor:normabs}
If $\Ot$ satisfies all hypotheses of Theorem~\ref{thm:NS} except that
$|\Ot|=|B_1|$ and $x_{\Ot}=0$ are replaced by
\[
\bigl|\,|\Ot|-|B_1|\,\bigr|+|x_{\Ot}|
\;\le\;\eta\,\|g\|_{L^2(\bdy B_1)},
\]
with $\eta\le\eta_0(N,\gamma_0)$ small enough, then there exist
$\widetilde{c}_*(N)>0$ and $\widetilde{\delta}_{\rm sph}(N,\gamma_0)>0$
such that conclusions~\eqref{eq:NSalpha}--\eqref{eq:NSfraenkel} hold with
$c_*(N)$ replaced by $\widetilde{c}_*(N)$. Equivalently, the
volume--barycenter normalisation in~\eqref{eq:normal} can be replaced by
an $O(\|g\|_{L^2(\bdy B_1)})$ defect at the cost of a dimensional
constant.
\end{corollary}

\begin{proof}[Reduction]
Form $\Ot'=\Ot-x_{\Ot}$ and $\Ot''=\lambda\,\Ot'$ with
$\lambda=(|B_1|/|\Ot'|)^{1/N}$. Then $|\lambda-1|\le C\eta\|g\|_{L^2}$
and $|x_{\Ot}|\le\eta\|g\|_{L^2}$. The new parametrisation $g''$ of
$\Ot''$ satisfies
$\|g''\|_{C^{2,\gamma_0}}\le\|g\|_{C^{2,\gamma_0}}
+C(N)\eta\|g\|_{L^2}\le 2\|g\|_{C^{2,\gamma_0}}$ for $\eta$ small. By
construction $|\Ot''|=|B_1|$ and $x_{\Ot''}=0$, and the perturbations of
$E,\alpha,\Asym$ are $O(\eta\|g\|_{L^2}^{2})$, absorbed at the cost of a
factor $\tfrac12$ in the conclusion. Setting
$\widetilde{c}_*(N)=c_*(N)/4$ closes the argument.
\end{proof}

\subsection{Sharpness}

The exponent $2$ on $\Asym$ in~\eqref{eq:NSfraenkel} (equivalently the
exponent $1/2$ on the deficit in the rearranged inequality
$\Asym(\Ot)\le C\,(E(\Ot)-E(B_1))^{1/2}$) cannot be improved, already
inside the nearly spherical class. This is shown in \cite[p.~2 and the
discussion preceding \S2]{BDV2013} by an ellipsoidal example.

\begin{remark}[Ellipsoidal sharpness]\label{rem:ellipsoid}
Let $E_\varepsilon=\{x:\sum x_i^{2}/a_i^{2}<1\}$ with $a_i=1+\varepsilon
e_i$, $\sum e_i=0$, $|e|\le 1$. A direct expansion gives
\[
\Asym(E_\varepsilon)=c_1(N,e)\,\varepsilon+O(\varepsilon^{2}),\qquad
E(E_\varepsilon)-E(B_1)=c_2(N,e)\,\varepsilon^{2}+O(\varepsilon^{3}),
\]
so $\Asym^{2}$ and the energy deficit are of the same order
$\varepsilon^{2}$. No inequality
$E(\Ot)-E(B_1)\ge c\,\Asym(\Ot)^{2-\theta}$ ($\theta>0$) can hold
uniformly even within nearly spherical sets, so the exponent $2$ on
$\Asym$ (equivalently $1/2$ on the deficit) is sharp.
\end{remark}

\section{Summary of constants}\label{sec:constants}

\begin{center}
\renewcommand{\arraystretch}{1.25}
\begin{tabular}{@{}lll@{}}
\toprule
\textbf{Symbol} & \textbf{Definition / origin} & \textbf{Where} \\
\midrule
$c_{\rm sph}(N)$ & $\dfrac{1}{32\,N^2}$, coefficient in BDV~Thm~3.3 &
\eqref{eq:BDV33}\\
$\delta_1(N,\gamma_0)$ & $C^{2,\gamma_0}$ smallness in BDV~Thm~3.3 &
Lem.~\ref{lem:BDV33}\\
$C_1(N)$ & $C_1\alpha\ge|\Omega\triangle B_1(x_\Omega)|^2$, BDV~4.2(i) &
\eqref{eq:BDV42i}\\
$\delta_3(N)$ & $L^\infty$ smallness in BDV~Lem.~4.2(iii) &
\eqref{eq:BDV42iii}\\
$C_3(N)$ & $\alpha\le C_3\|g\|_{L^2(\bdy B_1)}^2$, BDV~Lem.~4.2(iii) &
\eqref{eq:BDV42iii}\\
\midrule
$\delta_{\rm sph}(N,\gamma_0)$ & $\min(\delta_1,\delta_3)$ &
\eqref{eq:dsph}\\
$c_*(N)$ & $\dfrac{1}{32\,N^2\,C_3(N)}$ & \eqref{eq:cstar}\\
$c_*(N)|B_1|^{2}/C_1(N)$ & Fraenkel constant in~\eqref{eq:NSfraenkel} &
\eqref{eq:NSfraenkel}\\
$\widetilde{c}_*(N)$ & $c_*(N)/4$ after vol./barycenter normalisation &
Cor.~\ref{cor:normabs}\\
\bottomrule
\end{tabular}
\end{center}

\paragraph{Pure $N$-dependence.}
All constants depend only on $N$, except $\delta_1$ and hence
$\delta_{\rm sph}$ which depend additionally on $\gamma_0$. No radius
parameter $R$ enters, since by the time the closure is invoked the
upstream selection has returned a set whose convex hull is contained in
a ball of dimensional radius and whose volume is normalised to $|B_1|$.

\section{Use in Route~A}\label{sec:routeA}

In the Plan~1 / Route~A architecture, $\NS$ (Theorem~\ref{thm:NS}) is
applied as follows.

\begin{enumerate}
\item[(R1)] The upstream blocks (\texttt{SelectionPrinciple},
\texttt{GraphEntry}, \texttt{SchauderInterpolation}) produce a set
$\Ot$ which is $C^{2,\gamma_0}$ a graph over $\bdy B_1$ with
parametrisation $g$ satisfying $\|g\|_{C^{2,\gamma_0}}\le
h_{\rm sph}(N,\gamma_0,R)$ for some $h_{\rm sph}>0$.

\item[(R2)] After volume normalisation and barycenter shift
(Corollary~\ref{cor:normabs}) we may assume $|\Ot|=|B_1|$,
$x_{\Ot}=0$, and $h_{\rm sph}\le\delta_{\rm sph}(N,\gamma_0)$.

\item[(R3)] Theorem~\ref{thm:NS} then gives, with the explicit constants
of~\eqref{eq:cstar},
\[
E(\Ot)-E(B_1)
\;\ge\;\frac{|B_1|^{2}}{32\,N^2\,C_3(N)\,C_1(N)}\,\Asym(\Ot)^{2},
\]
the closing inequality of the sharp Saint--Venant stability with the
sharp exponent $2$ on $\Asym$ (equivalently $1/2$ on the deficit).
\end{enumerate}

\paragraph{What we do \emph{not} use.}
The proof of Theorem~\ref{thm:NS} does not use:
\begin{itemize}
\item BDV's selection principle (\cite[Prop.~4.4]{BDV2013});
\item BDV's compactness/contradiction step (\cite[Thm~4.3]{BDV2013});
\item Lemma~4.2(ii) (Lipschitz dependence of $\alpha$ on symmetric
difference).
\end{itemize}
The two BDV ingredients we do use (Theorem~3.3 and Lemma~4.2(i),(iii))
are both proved in \cite{BDV2013} by direct spherical-harmonic and
rearrangement computations.

\begin{thebibliography}{9}
\bibitem{BDV2013}
L.~Brasco, G.~De~Philippis, B.~Velichkov.
\emph{Faber--Krahn inequalities in sharp quantitative form}.
arXiv:1306.0392v1 (2013).
\end{thebibliography}

\end{document}
